\section{Introduction}
\label{sec:intro}
LLM-as-a-judge has emerged as a scalable alternative to human evaluation for assessing language model outputs~\citep{NEURIPS2023_91f18a12}. It enables practitioners to evaluate thousands of model outputs at a fraction of the cost and time required for human annotation~\citep{gu2024survey, tan2025judgebenchbench}. However, most existing work employs LLM judges for subjective evaluation~\citep{badshah-clev}, such as scoring outputs for helpfulness, harmlessness, or style in reward modeling, where no single correct answer exists and approximate agreement suffices~\citep{yuan-etal-2024-r, liu-etal-2023-g}. Extending LLM-as-a-judge to objective, factual evaluation remains an open challenge. 

In factual and open-domain QA, answers are either correct or incorrect, and the judge must make this determination reliably. However, obtaining ground-truth annotations for this purpose is expensive at scale~\citep{chianglee2023large, manas2024improving}, requires domain expertise for specialized questions, and is entirely infeasible when LLMs operate as autonomous agents answering open-ended queries in real-world environments where no pre-collected labels exist. LLM-as-a-judge offers a reference-free alternative that eliminates this dependency. Nevertheless, utilizing LLM judges to assess correctness without explicit references introduces two key reliability challenges~\citep{schroeder2025trustllm, badshah2025ref}.

First, LLM judges are bounded by their training data and may lack knowledge of recent events, rare entities, or specialized domains~\citep{10.1162/tacl_a_00754}. Such limitations can prevent reliable distinction between correct answers and plausible fabrications. Second, even if the judge has the relevant knowledge, hallucinations may still cause it to produce confident but incorrect verdicts~\citep{10.1145/3703155}. In both cases, the failures are silent as the judge provides no indication that its verdict is unreliable. Without a mechanism to quantify and act on this uncertainty, incorrect evaluations can propagate undetected into downstream decisions.

Selective evaluation provides a natural mitigation strategy by allowing the judge to abstain when uncertain and produce a verdict only when it is likely to be correct. Uncertainty quantification methods can identify unreliable verdicts~\citep{10.1145/3744238}. However, applying a heuristic threshold to these scores provides no formal guarantee on the error rate among accepted evaluations~\citep{scope}. Recent work on risk-controlled selective prediction addresses this gap by calibrating statistically valid thresholds that bound the False Discovery Rate (FDR) among selected outputs with high probability~\citep{safer, wang2025coin}. However, existing methods are designed to filter examinee answers in Question-Answering (QA) tasks rather than control errors in a meta-evaluation setting. Moreover, earlier methods treat abstention as the only recourse for uncertain instances. This discards cases where the judge's uncertainty stems from a knowledge gap or hallucination that could be resolved with additional evidence~\citep{tale_sb}.

In this paper, we propose a risk-controlled selective evaluation framework for LLM-as-a-judge that addresses the above limitations. Inspired by LLM-based agents that leverage external tools to overcome their parametric limitations, we equip the judge with web search as a tool to bridge its knowledge gaps. Our framework operates in two modes. In Mode~1, the judge model evaluates using only its parametric knowledge without any tool-augmentation. If the judge is confident, the verdict is accepted. If uncertain, the instance is routed to Mode~2 where the judge retrieves relevant evidence from the web and re-evaluates with this augmented context. If the judge remains uncertain after retrieval, it abstains and flags the instance for human review. Both thresholds are jointly calibrated on a held-out set with known ground-truth labels via Clopper--Pearson upper confidence bounds~\citep{be7c0fd0_cp}. We prove that the finite-sample FDR guarantee for a fixed threshold pair extends to the two-threshold routing setting without additional distributional assumptions. Our contributions are:
\begin{itemize}
\item We formulate risk-controlled selective evaluation for LLM-as-a-judge on factual QA.
\item We introduce a two-mode routing mechanism that equips the judge with adaptive web retrieval to resolve uncertainty before resorting to abstention.
\item We show that the Bernoulli structure required for the Clopper--Pearson bound extends to the two-threshold routing setting.
\item Experiments across open-domain QA benchmarks and judges of varying scales demonstrate that adaptive retrieval recovers substantial coverage over single-mode baselines while controlling the judge's error rate.
\end{itemize}